\documentclass[11pt,a4paper,oneside]{book}
\usepackage{90-faelles/latex/studiearkiv}
\begin{document}
\frontmatter
\begin{titlepage}
\centering
\vspace*{28mm}
{\sffamily\color{arkivblaa}\Large FYSIK · STUDIEARKIV\par}
\vspace{16mm}
{\Huge\bfseries Matematik til\\[5mm]universitetsfysik\par}
\vspace{12mm}
{\Large Håndbog, formler og forståelse\par}
\vfill
{\large Projektfundament og pilotkapitel\par}
\vspace{4mm}
{\large 26. oktober 2026\par}
\vspace{14mm}
\begin{minipage}{0.83\textwidth}
Denne første udgave indeholder et udviklet kapitel om differentialregning og en detaljeret disposition for de øvrige kapitler. Dispositionerne er markeret som planlagte og skal ikke læses som færdige noter. Indholdet er selvstændigt udarbejdet; der påstås ingen overensstemmelse med et bestemt KU-pensum.
\end{minipage}
\end{titlepage}
\chapter*{Læsevejledning}
\addcontentsline{toc}{chapter}{Læsevejledning}
Brug indholdsfortegnelsen og PDF-bogmærkerne til navigation. Differentialregningskapitlet samler definitioner, præcise forudsætninger, standardformler og metoder uden gennemregnede opgaveeksempler. Figurerne bruger dimensionsløse variable, medmindre andet er angivet. Fysiske variable skal altid fortolkes med deres enheder.

Notation: $\R$ betegner de reelle tal, $I$ et interval, og $f\colon I\to\R$ en reel funktion. $f'$ og $f^{(n)}$ betegner første og $n$-te afledte. $\dd$ er differentialtegnet, og $\e$ er eksponentialfunktionens grundtal. Antagelser står ved hvert resultat og gælder ikke automatisk på tværs af afsnit.

Kapitler udvikles trinvist efter repositoryets udviklingsvejledning. Formelsamlinger og konkrete opgaveløsninger til de enkelte fysikfag opbevares separat.
\begingroup\small
\tableofcontents
\endgroup
\mainmatter
\chapter{Algebra og funktioner}
\label{ch:algebra}
Algebra er sproget, der forbinder en matematisk model med dens ukendte størrelser. I fysik er opgaven sjældent blot at isolere et symbol: Man skal også holde styr på, hvilke værdier det må antage, hvilke omskrivninger der er reversible, og om det endelige udtryk har de rigtige dimensioner. Kapitlet behandler derfor algebra som et \emph{opslagsværktøj} for senere modeller og ikke blot som en samling regneregler.

\section{Tal og mængder}
\label{sec:tal-maengder}
En mængde fastlægger, hvilke objekter en matematisk påstand handler om.
Dette er lige så vigtigt som selve regneoperationen: Et heltalligt antal
partikler, en reel position og en kompleks amplitude er forskellige slags
størrelser. En ligning beskriver derfor først et veldefineret problem, når
både de tilladte værdier og betingelserne for dens udtryk er angivet.

\subsection{Talområder og deres regnemuligheder}
\label{sec:talomraader}
Vi bruger følgende konventioner:
\begin{center}
\begin{tabular}{@{}p{0.16\textwidth}p{0.35\textwidth}p{0.39\textwidth}@{}}
\toprule
Talområde & Definition eller beskrivelse & Typisk fysisk rolle\\
\midrule
$\mathbb N$ & $\{1,2,3,\ldots\}$ & Positive tælleindekser\\[3pt]
$\mathbb N_0$ & $\{0,1,2,\ldots\}$ & Antal partikler eller besættelsestal\\[3pt]
$\mathbb Z$ & $\{\ldots,-1,0,1,\ldots\}$ & Heltallige modeindekser\\[3pt]
$\mathbb Q$ & $\{p/q:p\in\mathbb Z,\ q\in\mathbb Z\setminus\{0\}\}$ & Rationale forhold og endelige decimaltal\\[3pt]
$\mathbb R$ & Den fuldstændige ordnede talkrop & Kontinuerte koordinater og parametre\\[3pt]
$\mathbb C$ & $\{a+ib:a,b\in\mathbb R,\ i^2=-1\}$ & Amplituder og svingningsrepræsentationer\\
\bottomrule
\end{tabular}
\end{center}
Andre bøger kan medregne nul i $\mathbb N$; her bruges $\mathbb N_0$, når
nul skal indgå. Med de sædvanlige identifikationer har vi kæden
\[
\boxed{\mathbb N\subsetneq\mathbb N_0\subsetneq\mathbb Z
\subsetneq\mathbb Q\subsetneq\mathbb R\subsetneq\mathbb C.}
\]
Symbolet $\subsetneq$ betyder en ægte delmængde: alle elementer i den
venstre mængde ligger i den højre, men de to mængder er ikke ens.
Vi bruger $\subseteq$, når lighed er tilladt.

Et talområde er \emph{lukket} under en operation, hvis operationens
resultat igen ligger i området, når den er defineret. De naturlige tal er
lukkede under addition og multiplikation, men ikke under vilkårlig
subtraktion. De hele tal tillader også subtraktion, mens division kan kræve
rationale tal. Både $\mathbb Q$, $\mathbb R$ og $\mathbb C$ er talkroppe:
de fire regningsarter kan udføres inden for området, bortset fra division
med nul. Kvadratroden af et positivt rationalt tal behøver derimod ikke
være rational, så regningsarternes lukkethed dækker ikke alle funktioner.

De reelle tal har desuden en fuldstændighedsegenskab: Enhver ikke-tom
opadtil begrænset delmængde af $\mathbb R$ har en mindste øvre grænse i
$\mathbb R$. Den egenskab gør det muligt at begrunde grænseværdier og
konvergens i bind 2. At $\mathbb Q$ ligger tæt i $\mathbb R$ betyder, at
ethvert åbent interval med forskellige endepunkter indeholder et rationalt
tal; det betyder ikke, at alle reelle tal er rationale.

Komplekse tal udvider løsningsmulighederne: Ligningen $x^2+1=0$ har ingen
reel løsning, men løsningerne $\pm i$ i $\mathbb C$. De komplekse tal
har imidlertid ingen total
ordensrelation, der gør dem til en ordnet talkrop. Hvis en sådan orden
fandtes, ville kvadrater af ikke-nul tal være positive; så skulle både
$1=1^2$ og $-1=i^2$ være positive, hvilket er umuligt. En ulighed som
$z>0$ bruges derfor ikke for et generelt komplekst tal. Man kan i stedet
undersøge dets reelle del eller dets reelle modulus, afhængigt af spørgsmålet.

\subsection{Mængdenotation og logiske udsagn}
\label{sec:maengdenotation}
En mængde er bestemt af sine elementer; deres rækkefølge og eventuelle
gentagelser i en opremsning ændrer ikke mængden. Notationen $x\in A$
betyder, at $x$ tilhører $A$, mens $x\notin A$ betyder det modsatte.
Den tomme mængde $\varnothing$ har ingen elementer. Bemærk forskellen på
$\varnothing$ og $\{0\}$: den sidste mængde har ét element.

Mængdebyggernotationen
\[
A=\{x\in U:P(x)\}
\]
betyder alle de elementer i grundmængden $U$, der opfylder betingelsen
$P(x)$. Det lodrette skel læses »for hvilke«. En grundmængde kan være et
talområde, et rum af vektorer eller senere et rum af funktioner.
\begin{definition}[Delmængde og lighed]
$A\subseteq B$ betyder, at hvert element i $A$ også ligger i $B$.
To mængder er ens netop når $A\subseteq B$ og $B\subseteq A$.
\end{definition}
At vise mængdelighed kræver således to inklusioner. Det er en generel
bevismetode, der senere bruges ved sammenligning af kerner, billeder og
løsningsrum i lineær algebra.

Kvantorerne $\forall$ og $\exists$ betyder henholdsvis »for alle« og
»der findes«. Deres rækkefølge er afgørende: I
$\forall x\in A\;\exists y\in B:P(x,y)$ må valget af $y$ afhænge af $x$;
i $\exists y\in B\;\forall x\in A:P(x,y)$ skal det samme $y$ virke for
alle $x$. Denne forskel bliver central for punktvis og uniform konvergens
i bind 2. Negation af en kvantificeret påstand følger reglerne
\[
\boxed{\begin{aligned}
\neg\bigl(\forall x\in A:P(x)\bigr)
&\ \Longleftrightarrow\ \exists x\in A:\neg P(x),\\
\neg\bigl(\exists x\in A:P(x)\bigr)
&\ \Longleftrightarrow\ \forall x\in A:\neg P(x).
\end{aligned}}
\]
Et enkelt modeksempel afkræfter altså en universel påstand. Udsagnet
$P\Rightarrow Q$ siger, at $P$ er en tilstrækkelig betingelse for $Q$,
og at $Q$ er en nødvendig betingelse for $P$. En ækvivalens
$P\Leftrightarrow Q$ kræver begge implikationer. Denne skelnen beskytter
mod at forveksle en nødvendig kontrol med en fuldstændig karakterisering.

\subsection{Mængdeoperationer som kombination af betingelser}
\label{sec:maengdeoperationer}
Antag, at $A,B\subseteq U$, hvor $U$ er en fast grundmængde. De vigtigste
operationer er samlet her:
\begin{center}
\begin{tabular}{@{}p{0.22\textwidth}p{0.39\textwidth}p{0.29\textwidth}@{}}
\toprule
Operation & Elementbetingelse & Betydning\\
\midrule
$A\cap B$ & $x\in A$ og $x\in B$ & Begge krav gælder\\[3pt]
$A\cup B$ & $x\in A$ eller $x\in B$ & Mindst ét krav gælder\\[3pt]
$A\setminus B$ & $x\in A$ og $x\notin B$ & Udeluk elementer fra $B$\\[3pt]
$A^c=U\setminus A$ & $x\in U$ og $x\notin A$ & Komplement relativt til $U$\\
\bottomrule
\end{tabular}
\end{center}
»Eller« er her inklusivt: et element i begge mængder ligger også i deres
foreningsmængde. Komplementet afhænger af $U$, og grundmængden skal derfor
fremgå af sammenhængen.

De vigtigste omskrivninger er distributivitet og De Morgans love:
\[
\boxed{\begin{aligned}
A\cap(B\cup C)&=(A\cap B)\cup(A\cap C),\\
A\cup(B\cap C)&=(A\cup B)\cap(A\cup C),\\
(A\cup B)^c&=A^c\cap B^c,\\
(A\cap B)^c&=A^c\cup B^c.
\end{aligned}}
\]
Disse regler følger af betingelserne på hvert enkelt element. For den
første komplementregel gælder for eksempel
\[
x\in(A\cup B)^c
\ \Longleftrightarrow\ x\notin A\ \text{og}\ x\notin B
\ \Longleftrightarrow\ x\in A^c\cap B^c.
\]
Det er en udledning af en generel identitet, ikke en regel, der blot skal
huskes fra en tegning.

For en indekseret familie $(A_j)_{j\in J}$ betyder
$\bigcup_{j\in J}A_j$, at mindst én betingelse opfyldes, mens
$\bigcap_{j\in J}A_j$ kræver alle betingelser opfyldt. Når mange fysiske
bindinger skal gælde samtidig, er snitmængden derfor den naturlige
beskrivelse af de tilladte tilstande.

\subsection{Intervaller, endepunkter og ubegrænsede områder}
\label{sec:intervaller}
For reelle $a<b$ defineres
\[
\begin{aligned}
[a,b]&=\{x\in\mathbb R:a\leq x\leq b\},&
(a,b)&=\{x\in\mathbb R:a<x<b\},\\
[a,b)&=\{x\in\mathbb R:a\leq x<b\},&
(a,b]&=\{x\in\mathbb R:a<x\leq b\}.
\end{aligned}
\]
Runde parenteser udelader endepunktet, firkantede parenteser medtager det.
I dansk litteratur ses også $]a,b[$ for $(a,b)$. Her bruges den sidstnævnte
notation konsekvent. Et interval indeholder alle reelle tal mellem to af
sine punkter; en foreningsmængde af adskilte intervaller behøver ikke selv
være et interval.

Symbolerne $+\infty$ og $-\infty$ er ikke reelle tal og medtages aldrig
som intervalendepunkter:
\[
[a,\infty)=\{x\in\mathbb R:x\geq a\},\qquad
(-\infty,b)=\{x\in\mathbb R:x<b\}.
\]
En ubegrænset mængde beskriver vilkårligt store eller små værdier; den
indeholder ikke et særligt tal »uendelig«. Hvis endepunkterne falder sammen,
er $[a,a]=\{a\}$, mens $(a,a)=\varnothing$.

Figur~\ref{fig:algebra-intervaller} viser forskellen mellem medtagne og
udeladte punkter og mellem et interval og et område med et hul.

\begin{figure}[htbp]
\centering
\begin{tikzpicture}[x=0.92cm,y=0.85cm,>=stealth,
  every node/.style={font=\small}]
\foreach \y in {0,1.3,2.6} {
  \draw[->,gray] (-0.4,\y)--(7,\y) node[right] {$x$};
}
\draw[very thick,arkivblaa] (1,2.6)--(5.6,2.6);
\draw[thick,draw=arkivblaa,fill=white] (1,2.6) circle (2.7pt);
\fill[arkivblaa] (5.6,2.6) circle (2.7pt);
\node[below] at (1,2.5) {$a$};
\node[below] at (5.6,2.5) {$b$};
\node[anchor=east] at (-0.6,2.6) {$(a,b]$};
\draw[very thick,arkivblaa] (1,1.3)--(3.2,1.3);
\draw[very thick,arkivblaa] (3.2,1.3)--(5.6,1.3);
\foreach \x in {1,5.6} {
  \draw[thick,draw=arkivblaa,fill=white] (\x,1.3) circle (2.7pt);
}
\draw[thick,draw=arkivrod,fill=white] (3.2,1.3) circle (2.7pt);
\node[below] at (1,1.2) {$a$};
\node[below] at (3.2,1.2) {$c$};
\node[below] at (5.6,1.2) {$b$};
\node[anchor=east] at (-0.6,1.3) {$(a,b)\setminus\{c\}$};
\draw[very thick,arkivblaa] (1.4,0)--(5.2,0);
\foreach \x in {1.4,5.2} {
  \draw[thick,draw=arkivblaa,fill=white] (\x,0) circle (2.7pt);
}
\draw (3.3,-0.08)--(3.3,0.08);
\node[below] at (3.3,-0.1) {$x_0$};
\draw[<->,arkivrod] (3.3,0.4)--(5.2,0.4)
  node[midway,above] {$r$};
\node[below] at (1.4,-0.1) {$x_0-r$};
\node[below] at (5.2,-0.1) {$x_0+r$};
\node[anchor=east] at (-0.6,0) {$|x-x_0|<r$};
\end{tikzpicture}
\caption{Intervaller på en skematisk, dimensionsløs talakse. En fyldt prik
medtager et endepunkt, mens en åben prik udelader det. I midten antages
$a<c<b$: udeladelsen af $c$ deler området i to intervaller. Nederst vises
afstandsbetingelsen omkring $x_0$ med $r>0$.}
\label{fig:algebra-intervaller}
\end{figure}


Afstand på talaksen skrives $d(x,y)=|x-y|$. For en radius $r>0$ kan et
åbent interval derfor også beskrives som
\[
\boxed{\{x\in\mathbb R:|x-x_0|<r\}=(x_0-r,x_0+r).}
\]
Udledningen følger af $|u|<r\Leftrightarrow-r<u<r$ med $u=x-x_0$.
For $|x-x_0|\leq r$ medtages begge endepunkter. Denne notation beskriver
både tolerancer omkring en referenceværdi og de små omgivelser, der bruges
ved kontinuitet i bind 2. Absolutværdi og fejlgrænser behandles videre i
afsnittet om uligheder i dette kapitel.

\subsection{Produktmængder og flere variable}
\label{sec:produktmaengder}
Produktmængden $A\times B$ består af ordnede par:
\[
A\times B=\{(a,b):a\in A,\ b\in B\}.
\]
Rækkefølgen er en del af informationen. Generelt er $(a,b)$ ikke det
samme som $(b,a)$, selv om mængderne $\{a,b\}$ og $\{b,a\}$ er ens.
Notationen $\mathbb R^n$ betyder produktet af $n$ kopier af $\mathbb R$;
et element er en koordinatliste $(x_1,\ldots,x_n)$.

En produktmængde tillader uafhængige valg fra hver faktor. En binding
$F(x,y)=0$ udvælger derimod en delmængde af $\mathbb R^2$, hvor
koordinaterne ikke kan vælges frit. Dette er forudsætningen for senere at
beskrive løsningsmængder i lineær algebra, niveaukurver i bind 2 og
geometriske bindinger i bind 5. En koordinatliste alene fastlægger ikke et
fysisk vektorrum med enheder, norm og basis; den ekstra struktur skal
angives, når den bruges.

\subsection{Definitionsmængder og reversible omskrivninger}
\label{sec:definitionsmaengder}
Hvis et udtryk $F(x)$ skal være reelt, findes dets naturlige
\emph{definitionsmængde} ved at kræve, at alle dets dele er definerede.
Løsningsmængden for en ligning er derefter
\[
\boxed{S=\{x\in D:F(x)=0\},}
\]
hvor $D$ er den relevante definitionsmængde efter eventuelle yderligere
modelkrav. Betingelser skal kombineres med snitmængde, når alle dele af
udtrykket skal være definerede samtidigt:
\begin{center}
\begin{tabular}{@{}p{0.33\textwidth}p{0.57\textwidth}@{}}
\toprule
Reelt udtryk & Nødvendig definitionsbetingelse\\
\midrule
$P(x)/Q(x)$ & $Q(x)\ne0$\\[3pt]
$\sqrt{g(x)}$ & $g(x)\geq0$\\[3pt]
$1/\sqrt{g(x)}$ & $g(x)>0$\\[3pt]
$\ln g(x)$ & $g(x)>0$\\
\bottomrule
\end{tabular}
\end{center}
Desuden skal de indre udtryk være definerede. En bogstavelig formel kan
have et snævrere domæne end en forkortet repræsentation. Forkortning
\[
\frac{(x-a)P(x)}{x-a}=P(x)\qquad (x\ne a)
\]
beviser kun lighed på det oprindelige domæne. At tildele venstresiden en
værdi ved $x=a$ er en udvidelse af funktionen, ikke blot en omskrivning.

Division med et udtryk, der kan være nul, kan fjerne løsninger, mens
multiplikation med det kan tilføje løsninger. Kvadrering er heller ikke
generelt reversibel. For reelle $x,a$ gælder
\[
x=\sqrt a\quad\Longleftrightarrow\quad
\bigl(x^2=a\ \text{og}\ x\geq0\bigr)\qquad(a\geq0).
\]
Uden fortegnskravet på $x$ har man kun implikationen fra venstre mod
højre. Kandidater, der fremkommer ved en sådan omskrivning, skal derfor
kontrolleres i den oprindelige ligning. Det samme princip gælder, når
parameterafhængige nævnere eller inversgrene indgår.

\subsection{Metodevalg og fysisk gyldighed}
\label{sec:maengder-fysik}
Ved et algebraisk modelproblem er en nyttig rækkefølge at angive
grundmængden, opstille definitionsbetingelserne, skære dem med de fysiske
krav og først derefter løse ligningen. Ved flere tilladte alternativer
bruges en foreningsmængde. En parameter, hvis fortegn påvirker kravene,
skal behandles særskilt; et formelt resultat for positive parametre kan
ikke uden kontrol anvendes ved nul eller negative værdier.

For en endimensional konservativ bevægelse med konstant masse $m>0$ og
potentiel energi $V(q)$ giver energirelationen
\[
E=\frac{p^2}{2m}+V(q),\qquad
p^2=2m\bigl(E-V(q)\bigr).
\]
Et reelt momentum kræver $V(q)\leq E$. De energitilladte positioner er
således mængden
\[
A_E=\{q\in D_V:V(q)\leq E\}.
\]
Den kan bestå af flere adskilte intervaller; en partikel følger ikke
automatisk alle disse intervaller fra én begyndelsestilstand.
Ved et regulært vendepunkt med $V(q)=E$ er $p=0$. For $V(q)<E$ giver
$p=\pm\sqrt{2m(E-V(q))}$ to momentumgrene. Energiens nødvendige
positionsbetingelse vælger altså hverken bevægelsesretningen eller en
fuld tidsudvikling. Denne forskel mellem en tilladt mængde og en faktisk
løsning er vigtig i mekanik og i bind 3 om differentialligninger.

I statistisk fysik kan et besættelsestal ligge i $\mathbb N_0$, mens en
klassisk koordinat beskrives i $\mathbb R$. En kontinuert approksimation
til et diskret antal er en modelbeslutning og kræver sin egen begrundelse.
Mængdenotation holder disse valg synlige uden at foregribe den senere
fysiske teori.

\paragraph{Forbindelser i værket.}
Funktionsafsnittet nedenfor bruger domæne, codomæne og billede til at
præcisere, hvornår en invers eksisterer. Kapitlet om lineær algebra i dette
bind tilføjer regneoperationer og lineære bindinger på vektormængder.
Bind 2 bruger intervaller og kvantorer til grænseværdier, og bind 5
beskriver områder og koordinater i flere dimensioner. Disse henvisninger
angiver emneplacering; de er ikke en påstand om, at alle senere afsnit er
færdigskrevet.

\section{Potenser, rødder og brøker}
\label{sec:potenser-broeker}
Potenser samler gentagen multiplikation i ét symbol. Rødder vender denne
operation om, og brøker gør division til en del af regningen. De tre
værktøjer hænger tæt sammen, men deres regneregler har forskellige
forudsætninger. Særligt fortegn, nulværdier og ikke-heltallige eksponenter
kræver opmærksomhed, før en formel kan forkortes.

\subsection{Heltallige potenser og nulværdier}
For $n\in\mathbb N$ er $a^n$ produktet af $n$ faktorer lig med $a$.
For $a\ne0$ udvides definitionen ved
\[
a^0=1,\qquad a^{-n}=\frac{1}{a^n}.
\]
Med disse definitioner gælder for $m,n\in\mathbb Z$ og $a\ne0$
\[
\boxed{a^m a^n=a^{m+n},\qquad
\frac{a^m}{a^n}=a^{m-n},\qquad (a^m)^n=a^{mn}.}
\]
For positive eksponenter følger den første regel ved at tælle antallet af
faktorer. Kravet om at bevare samme regel ved $n=-m$ giver
$a^m a^{-m}=a^0=1$ og begrunder definitionen af negative potenser.
For $a=0$ er positive heltallige potenser definerede, men negative potenser
og $0^0$ bruges ikke her. En forkortning må derfor ikke skjule et
uoprindeligt defineret nuludtryk.

For heltallige $n$ gælder også
\[
(ab)^n=a^n b^n,\qquad
\left(\frac ab\right)^n=\frac{a^n}{b^n},
\]
hvor alle nævnere og negative potenser skal være definerede. En potens
fordeler sig over et produkt, men ikke generelt over en sum.
Fortegn og parenteser skal desuden holdes adskilt:
$(-a)^2=a^2$, mens $-a^2=-(a^2)$.

\subsection{Reelle rødder og rationale eksponenter}
For $n\in\mathbb N$ defineres den reelle $n$'te rod ved
\[
\begin{cases}
\sqrt[n]{a}=\text{det entydige }r\geq0\text{ med }r^n=a,
 & n\text{ lige},\ a\geq0,\\
\sqrt[n]{a}=\text{det entydige reelle }r\text{ med }r^n=a,
 & n\text{ ulige},\ a\in\mathbb R.
\end{cases}
\]
En lige rod er altså valgt ikke-negativ. Ligningen $x^2=a$ og selve
symbolet $\sqrt a$ har forskellige roller: for $a>0$ har ligningen to
rødder, men symbolet betegner kun den positive. Derfor gælder
\[
\boxed{\sqrt{x^2}=|x|\quad(x\in\mathbb R),\qquad
(\sqrt x)^2=x\quad(x\geq0).}
\]
Den første identitet følger, fordi $|x|$ er ikke-negativ og har kvadrat
$x^2$. At fjerne absolutværdien kræver den ekstra betingelse $x\geq0$.

For $a>0$ og heltal $p$, $q\in\mathbb N$ sætter man
\[
a^{p/q}=\bigl(\sqrt[q]{a}\bigr)^p.
\]
Værdien afhænger ikke af valget af brøkrepræsentation: Hvis $p/q=r/s$,
er de to positive kandidaters $qs$'te potenser begge $a^{ps}=a^{rq}$;
entydigheden af den positive rod giver samme værdi. Her må $p$ og $r$
også være negative, da $a\ne0$.
For $a=0$ kan positive rationale potenser defineres som nul.

For negative baser bruges en mere begrænset reel konvention:
Er $p/q$ forkortet med $q>0$ ulige, kan man sætte
$a^{p/q}=(\sqrt[q]{a})^p$. Ved lige nævner findes ingen reel værdi.
Man skal forkorte eksponenten først; ellers kan et valg af rod i en
mellemregning give et andet resultat. De generelle regler for sammensatte
potenser kan ikke overføres uden kontrol: Hvis $a<0$, er
$(a^2)^{1/2}=|a|\ne a^{2(1/2)}$.

\subsection{Vilkårlige reelle eksponenter og regneoversigt}
For en positiv base defineres vilkårlige reelle potenser ved
\[
\boxed{a^\alpha=\exp(\alpha\ln a)\qquad(a>0,\ \alpha\in\mathbb R).}
\]
Denne definition udvider de rationale potenser og giver de almindelige
potensregler uden grenproblemer. Reglerne følger af
$\exp(u+v)=\exp u\exp v$ og $\ln(ab)=\ln a+\ln b$ for positive
argumenter; eksponential- og logaritmefunktionerne behandles samlet
senere i dette kapitel.
\begin{center}
\begin{tabular}{@{}p{0.56\textwidth}p{0.34\textwidth}@{}}
\toprule
Regel & Tilstrækkelige reelle betingelser\\
\midrule
$a^\alpha a^\beta=a^{\alpha+\beta}$ & $a>0$\\[3pt]
$a^\alpha/a^\beta=a^{\alpha-\beta}$ & $a>0$\\[3pt]
$(a^\alpha)^\beta=a^{\alpha\beta}$ & $a>0$\\[3pt]
$(ab)^\alpha=a^\alpha b^\alpha$ & $a,b>0$\\[3pt]
$(a/b)^\alpha=a^\alpha/b^\alpha$ & $a,b>0$\\[3pt]
$\sqrt{ab}=\sqrt a\sqrt b$ & $a,b\geq0$\\[3pt]
$\sqrt{a/b}=\sqrt a/\sqrt b$ & $a\geq0,\ b>0$\\
\bottomrule
\end{tabular}
\end{center}
Betingelserne i tabellen er sikre fælles domæner; enkelte identiteter kan
være gyldige på større områder ved bestemte heltallige eksponenter.
Den reelle definition giver ikke vilkårlige irrationale potenser af
negative tal. Komplekse potenser kræver logaritmegrene og hører til bind 4.

\subsection{Brøkregning og fællesnævnere}
Brøken $a/b$ er for $b\ne0$ det tal, som multipliceret med $b$ giver $a$.
For reelle tal gælder
\[
\boxed{\begin{aligned}
\frac ab\pm\frac cd&=\frac{ad\pm bc}{bd}
&& (b,d\ne0),\\
\frac ab\,\frac cd&=\frac{ac}{bd}
&& (b,d\ne0),\\
\frac{a/b}{c/d}&=\frac{ad}{bc}
&& (b,c,d\ne0).
\end{aligned}}
\]
Additionsreglen udledes ved at skrive $a/b=ad/(bd)$ og
$c/d=bc/(bd)$. Man omskriver altså til samme enhed, før tællerne lægges
sammen. Divisionsreglen følger ved at gange med det reciprokke af den
ikke-nul divisor $c/d$.

Forkortning er division af hele tælleren og hele nævneren med en fælles
\emph{faktor}. Man kan ikke forkorte et enkelt led gennem en sum.
Hvis en fælles faktor afhænger af en variabel, skal dens nulpunkter
fortsat udelukkes fra det oprindelige domæne, jf.
afsnit~\ref{sec:definitionsmaengder}. Ved rationelle funktioner vælges en
fællesnævner efter faktorerne og deres nødvendige potenser; produktet af
alle nævnere virker, men kan være unødigt stort.

\subsection{Rationalisering og metodevalg}
Et rodudtryk i en forskel kan omskrives med den konjugerede faktor.
For $u,v\geq0$ og $u+v>0$ gælder
\[
\sqrt u-\sqrt v=\frac{u-v}{\sqrt u+\sqrt v}.
\]
Identiteten kommer af
$(\sqrt u-\sqrt v)(\sqrt u+\sqrt v)=u-v$.
Omvendt kan nævneren $\sqrt u-\sqrt v$ rationaliseres, når $u\ne v$:
\[
\frac{1}{\sqrt u-\sqrt v}
=\frac{\sqrt u+\sqrt v}{u-v}.
\]
Omskrivningerne er nyttige ved grænseværdier og ved forskelle mellem
næsten lige rødder. De udvider ikke automatisk domænet. Ved numerisk
beregning kan den første form reducere tabet af cifre i selve
rodsubtraktionen, men den samlede nøjagtighed afhænger også af, hvordan
$u-v$ beregnes; fejlanalyse hører til bind 4.

Vælg heltallige potensregler ved gentagne faktorer og reciprokke udtryk,
roddefinitionen ved spørgsmål om fortegn og reelle løsninger og positive
baser, når vilkårlige reelle eksponenter skal samles. Ved summer af brøker
findes fællesnævneren før forkortning. Ved rodligninger kontrolleres de
oprindelige fortegnskrav efter enhver kvadrering.

\subsection{Dimensionskontrol og fysiske skaleringslove}
Hvis $X$ er en fysisk størrelse med dimension $[X]$, har en heltallig
potens dimension $[X^n]=[X]^n$. Rødder og rationale potenser giver
tilsvarende brøkpotenser af enhederne, når størrelsen og den valgte
rodgren er meningsfulde. Addition kræver samme dimension i alle led,
og en algebraisk lighed kræver samme dimension på begge sider.

For en dimensionsbærende positiv størrelse er en skalering med en
vilkårlig eksponent klarest skrevet ved en positiv referenceskala $X_0$:
\[
\boxed{Y=Y_0\left(\frac{X}{X_0}\right)^\alpha.}
\]
Brøken $X/X_0$ er dimensionsløs, mens $Y_0$ har samme dimension som $Y$.
En potenslov $Y=CX^\alpha$ kan også bruges, men da afhænger dimensionen
af $C$ af $\alpha$. En variabel eksponent bør derfor ikke sættes på en
dimensionsbærende størrelse uden referenceskala.

En elastisk energi $U=kq^2/2$ med $k>0$ giver to mulige fortegn for
forskydningen, når energien og $k$ kendes. Kvadratroden bestemmer
størrelsen $|q|=\sqrt{2U/k}$, ikke bevægelsesretningen. Med
$[k]=\text{energi}/\text{længde}^2$ får roden dimension længde.
Dette forbinder rodens grenvalg med de tilladte tilstande fra
afsnit~\ref{sec:maengder-fysik}; den faktiske bevægelse kræver dynamik i
bind 3 og mekaniknoterne.

\section{Polynomier, faktorisering og binomialformlen}
\label{sec:polynomier}
Et polynomium bygger en funktion op af endeligt mange heltallige potenser.
Denne enkle form gør det muligt at samle led, finde nulpunkter og
omskrive en model uden at ændre dens indhold. Polynomier beskriver både
præcise algebraiske relationer og lokale tilnærmelser til mere komplicerede
funktioner. En faktorisering fremhæver rødderne, mens den udfoldede form
fremhæver koefficienterne; den nyttigste form afhænger af spørgsmålet.

\subsection{Koefficienter, grad og polynomiumsidentitet}
Et polynomium over $\mathbb K=\mathbb R$ eller $\mathbb C$ skrives
\[
P(x)=\sum_{k=0}^n a_k x^k=a_0+a_1x+\cdots+a_nx^n,
\qquad a_k\in\mathbb K.
\]
Hvis $a_n\ne0$, er graden $\deg P=n$, og $a_n$ er den ledende
koefficient. Et ikke-nul konstant polynomium har grad nul; nulpolynomiet
tilordnes ikke en endelig grad her. For ikke-nul polynomier gælder
\[
\deg(PQ)=\deg P+\deg Q.
\]
Dette følger, fordi produktets højeste led er
$a_n b_m x^{n+m}$ med $a_nb_m\ne0$. Derimod kan led annulleres ved
addition, så $\deg(P+Q)\leq\max(\deg P,\deg Q)$, hvis summen er ikke-nul.

To polynomier er identiske, når koefficienterne for hver potens er ens.
En identitet gælder for alle $x$ i talområdet; en ligning kan være sand
kun for bestemte værdier. Over $\mathbb R$ eller $\mathbb C$ giver lighed
som funktioner også lighed af koefficienter: et ikke-nul polynomium af
grad $n$ kan højst have $n$ forskellige rødder, som begrundes nedenfor.
Man kan derfor sammenligne koefficienter, når en omskrivning skal gælde
identisk. Et tjek i blot nogle få punkter beviser derimod ikke generelt
en identitet uden kontrol af graden.

\subsection{Standardfaktoriseringer og deres udledning}
Distributiviteten $a(b+c)=ab+ac$ giver følgende opslag:
\[
\boxed{\begin{aligned}
(a+b)^2&=a^2+2ab+b^2,\\
(a-b)^2&=a^2-2ab+b^2,\\
a^2-b^2&=(a-b)(a+b),\\
(a+b)^3&=a^3+3a^2b+3ab^2+b^3,\\
(a-b)^3&=a^3-3a^2b+3ab^2-b^3,\\
a^3-b^3&=(a-b)(a^2+ab+b^2),\\
a^3+b^3&=(a+b)(a^2-ab+b^2).
\end{aligned}}
\]
Kvadratsætningen fås ved at gange to faktorer $(a+b)$ ud; de to blandede
led er begge $ab$. Kubiksætningen følger ved at gange resultatet endnu
en gang med $(a+b)$. Det er netop de blandede led, der gør en potens af
en sum forskellig fra summen af potenser.

En generel forskel af potenser faktoriseres som
\[
a^n-b^n=(a-b)\sum_{k=0}^{n-1}a^{n-1-k}b^k
\qquad(n\in\mathbb N).
\]
Ved udmultiplikation annulleres alle mellemliggende led parvist.
Hvis $a\ne b$, kan man dividere med $a-b$; ved $a=b$ gælder den
uforkortede identitet stadig. Med $a=r$, $b=1$ giver samme idé
\[
\sum_{k=0}^{n-1}r^k=\frac{1-r^n}{1-r}\quad(r\ne1),
\qquad \sum_{k=0}^{n-1}1=n.
\]
Dette er en \emph{endelig} sum og kræver ingen konvergensbetingelse.
Uendelige geometriske rækker kræver en særskilt grænseovergang i bind 2.

\subsection{Binomialkoefficienter og binomialformlen}
Når produktet af $n$ faktorer $(a+b)$ udfoldes, opstår et led
$a^{n-k}b^k$ ved at vælge $b$ fra præcis $k$ af faktorerne. Antallet af
sådanne valg er binomialkoefficienten
\[
\binom nk=\frac{n!}{k!(n-k)!},\qquad
n\in\mathbb N_0,\quad k=0,\ldots,n,\quad 0!=1.
\]
Der er $n(n-1)\cdots(n-k+1)$ ordnede valg af $k$ forskellige positioner;
division med $k!$ fjerner rækkefølgen af de valgte positioner.
Dermed fås binomialformlen
\[
\boxed{(a+b)^n=\sum_{k=0}^n\binom nk a^{n-k}b^k
\qquad(n\in\mathbb N_0).}
\]
Koefficienterne er symmetriske, $\binom nk=\binom n{n-k}$, og opfylder
Pascals relation
\[
\binom nk=\binom{n-1}{k-1}+\binom{n-1}{k}\qquad(1\leq k\leq n-1).
\]
Et valg indeholder enten den sidste position eller gør det ikke; det
opdeler valgene i netop de to grupper. For $n=0$ forstås produktet som
et tomt produkt lig med 1; dette er ikke en generel definition af $0^0$
for vilkårlige reelle potenser.

Formlen gælder for almindelige reelle eller komplekse tal, hvor
multiplikationen er kommutativ. Den må ikke uden videre anvendes på
matricer: $(A+B)^2=A^2+AB+BA+B^2$ reduceres kun til den almindelige
kvadratsætning, hvis $AB=BA$. Potenser med ikke-heltallig eksponent giver
ikke denne endelige sum; binomialrækken behandles med sine
konvergensbetingelser i bind 2.

\subsection{Polynomiumsdivision, restsætning og rødder}
For polynomier $P,Q$ over $\mathbb R$ eller $\mathbb C$ med $Q\ne0$
findes entydige polynomier $S,R$, så
\[
\boxed{P(x)=Q(x)S(x)+R(x),\qquad R=0\ \text{eller}\ \deg R<\deg Q.}
\]
Metoden er at dividere det aktuelt højeste led i resten med det ledende
led i $Q$, føje kvotienten til $S$ og trække det tilsvarende multiplum af
$Q$ fra. Graden af den tilbageværende rest falder ved hvert trin, så
processen stopper. Hvis der fandtes to fremstillinger, ville
$Q(S_1-S_2)=R_2-R_1$ have grad mindst $\deg Q$ på venstresiden og mindre
på højresiden, medmindre begge sider er nul. Det begrunder entydigheden.

Ved division med $x-a$ er resten konstant. Indsættelse af $x=a$ giver
restsætningen og faktorsætningen:
\[
\boxed{P(x)=(x-a)S(x)+P(a),\qquad
P(a)=0\ \Longleftrightarrow\ (x-a)\text{ er faktor i }P.}
\]
Et ikke-nul polynomium af grad $n$ har derfor højst $n$ forskellige
rødder: Hver forskellig rod fjerner en lineær faktor, og graden falder
én ad gangen. En rod $a$ har multiplicitet $m$, hvis
$P(x)=(x-a)^m H(x)$ med $H(a)\ne0$. Multipliciteten er altså antallet af
gange faktoren forekommer, ikke antallet af forskellige rodværdier.

Algebraens fundamentalsætning siger, at et polynomium af grad $n\geq1$
over $\mathbb C$ kan faktoriseres i $n$ lineære faktorer talt med
multiplicitet. Sætningen anføres her uden bevis. Har polynomiet reelle
koefficienter, gælder $P(\bar z)=\overline{P(z)}$, så ikke-reelle rødder
kommer i konjugerede par. Et par $a\pm ib$ med $b\ne0$ giver faktoren
\[
\bigl(x-(a+ib)\bigr)\bigl(x-(a-ib)\bigr)
=(x-a)^2+b^2,
\]
som er irreducibel over $\mathbb R$. Et reelt polynomium kan derfor
faktoriseres i reelle lineære og irreducible kvadratiske faktorer.
Det er grundlaget for partialbrøksopdelingen senere i kapitlet.

\subsection{Metodevalg, fortegn og fysiske forbindelser}
Faktorisér, når nulpunkter, fælles faktorer eller fortegn er det centrale;
udfold og sammenlign koefficienter, når en identitet eller en
parameterbestemmelse skal kontrolleres. Brug polynomiumsdivision, når en
brøks tæller har mindst samme grad som nævneren, eller når et kendt
nulpunkt skal skilles ud. En eksakt faktoriseringsmulighed er nyttig,
men der findes ikke en generel rodformel med radikaler for alle
polynomier af grad fem eller højere; numeriske rødder behandles i bind 4.

Omkring en reel rod $a$ med multiplicitet $m$ er fortegnet af
$P(x)=(x-a)^m H(x)$ bestemt af $x-a$ og det lokalt uændrede fortegn af
$H$. Derfor skifter $P$ fortegn gennem roden ved ulige $m$, men ikke ved
lige $m$. Denne sammenhæng mellem faktorisering og grafen er vist i
figur~\ref{fig:algebra-multiplicitet} og bruges ved uligheder og ved
undersøgelse af, hvor en fysisk størrelse kan være positiv.

\begin{figure}[htbp]
\centering
\begin{tikzpicture}
\begin{axis}[width=0.79\textwidth,height=6cm,
 axis lines=middle,xmin=-1.15,xmax=1.15,ymin=-1.15,ymax=1.3,
 xlabel={$u=x-a$},ylabel={$P$},
 xtick={-1,0,1},ytick={-1,1},
 samples=120,domain=-1.05:1.05,
 legend style={at={(0.03,0.97)},anchor=north west,
 draw=none,fill=white,font=\small},
 label style={font=\small}]
\addplot[arkivblaa,very thick] {x^2};
\addlegendentry{$P=u^2$: lige multiplicitet}
\addplot[arkivrod,very thick,dashed] {x^3};
\addlegendentry{$P=u^3$: ulige multiplicitet}
\end{axis}
\end{tikzpicture}
\caption{Dimensionsløse grundformer ved en rod $x=a$.
En rod med lige multiplicitet berører nul uden fortegnsskift;
en rod med ulige multiplicitet krydser nul. For et generelt polynomium
$(x-a)^mH(x)$ gælder fortegnsargumentet lokalt, hvor $H$ er forskellig
fra nul; $H$ kan ændre både skala og kurvens orientering.}
\label{fig:algebra-multiplicitet}
\end{figure}


For en dimensionsbærende koordinat $q$ skal koefficienterne i en energi
$V(q)=\sum_k c_k q^k$ opfylde $[c_k]=[\text{energi}]/[q]^k$.
Ved symmetrien $V(-q)=V(q)$ for alle $q$ må de ulige koefficienter være
nul; dette følger ved at sammenligne koefficienter. Et lige potentiale
kan derfor organiseres med lige potenser uden at gætte dets fysiske
parametre. En lokal polynomiumsapproksimation er derimod ikke en
identitet for hele potentialet; dens restled og gyldighedsinterval
behandles i bind 2 om Taylorudvikling.

I lineær algebra er $\det(\lambda I-A)$ et polynomium i $\lambda$.
Faktorisering og multiplicitet vender derfor tilbage ved egenværdier,
men en gentagen polynomiumsrod garanterer ikke flere uafhængige
egenvektorer. Den ekstra geometriske betingelse behandles i kapitlet om
lineær algebra i dette bind.

\section{Ligninger og uligheder}
\subsection{Andengradsligninger og kvadratkomplettering}
Den generelle andengradsligning $ax^2+bx+c=0$ med $a\ne0$ har rødder
\[
\boxed{x_\pm=\frac{-b\pm\sqrt{b^2-4ac}}{2a}}.
\]
Diskriminanten $\Delta=b^2-4ac$ afgør antallet af reelle rødder. Fuldstændiggørelse af kvadratet giver
\[
ax^2+bx+c=a\left(x+\frac{b}{2a}\right)^2+
c-\frac{b^2}{4a}.
\]
Denne form er ofte mere informativ end rodformlen, når et minimum af en energifunktion eller en kvadratisk approksimation skal bestemmes.

\subsection{Brøker og partialbrøker}
En rational funktion $P(x)/Q(x)$ kan ofte dekomponeres efter polerne i nævneren. For $a\ne b$ er eksempelvis
\[
\frac{1}{(x-a)(x-b)}
=\frac{1}{a-b}\left(\frac{1}{x-a}-\frac{1}{x-b}\right).
\]
Gentagne lineære faktorer $(x-a)^m$ kræver led med alle nævnerpotenser $(x-a)^{-k}$ for $1\leq k\leq m$. Irreducible reelle andengradsfaktorer kræver en lineær tæller. Hvis tællerens grad ikke er mindre end nævnerens, udføres først polynomiumsdivision. Metoden bruges især før integration eller Laplace-inversion.

\subsection{Uligheder og absolutværdi}
Absolutværdi er afstanden fra nul på den reelle talakse:
\[
|x|=\begin{cases}x,&x\geq0,\\-x,&x<0.\end{cases}
\qquad |x+y|\leq |x|+|y|.
\]
Hvis en ulighed multipliceres med et negativt tal, skifter ulighedstegnet retning. Ved multiplikation med en ukendt parameter skal man derfor skelne mellem dens fortegnstilfælde.

For $a>0$ gælder den nyttige ækvivalens
\[
|x-x_0|<a\quad\Longleftrightarrow\quad x_0-a<x<x_0+a.
\]
I perturbationsteori giver en antagelse som $|\varepsilon|\ll1$ ikke alene en streng matematisk fejlgrænse. En kontrolleret tilnærmelse skal kobles til et restled eller en eksplicit størrelsesorden.

\section{Funktioner, inverse funktioner og sammensætning}
\subsection{Domæne, billede og sammensætning}
En funktion $f:D\to Y$ tildeler hvert $x\in D$ præcis én værdi $f(x)\in Y$. Definitionsmængden $D$, codomænet $Y$ og billedmængden $f(D)$ er forskellige begreber. En fysisk formel kan være algebraisk meningsfuld for et større interval end det, hvor modellen faktisk gælder.

Sammensætningen $f\circ g$ defineres ved $(f\circ g)(x)=f(g(x))$, når $x$ ligger i $g$'s definitionsmængde, og $g(x)$ desuden ligger i $f$'s definitionsmængde. Domænet for den sammensatte funktion er derfor
\[
D_{f\circ g}=\{x\in D_g:g(x)\in D_f\}.
\]
Et vigtigt modelspørgsmål er, om en funktion er injektiv (forskellige input giver forskellige output), surjektiv (hele codomænet rammes) eller bijektiv (begge dele). Kun en bijektiv funktion har en entydig global invers som afbildning mellem de angivne mængder.

\subsection{Inverse funktioner og implicitte relationer}
Hvis $f$ er strengt monoton på et interval, er den injektiv dér og har en invers $f^{-1}:f(D)\to D$. Hvis $f$ er differentiabel, og $f'(x)\ne0$, gælder lokalt
\[
\boxed{(f^{-1})'(f(x))=\frac{1}{f'(x)}}.
\]
Antagelsen $f'(x)\ne0$ er afgørende. For eksempel har $f(x)=x^3$ en global invers, men formlen giver ingen endelig afledt for inversen ved $f(0)=0$.

En implicit relation $F(x,y)=0$ beskriver ikke nødvendigvis en global funktion $y(x)$. Hvis $F$ er kontinuert differentiabel og $F_y(x_0,y_0)\ne0$, garanterer den implicitte funktionssætning en lokal gren med
\[
\frac{\dd y}{\dd x}=-\frac{F_x}{F_y}.
\]
Det er et typisk værktøj, når en bane eller en bevaret energi definerer en relation mellem størrelser, som er besværlig at løse eksplicit.

\subsection{Skalering, dimensionsløse variable og følsomhed}
En størrelse $x$ med dimension kan gøres dimensionsløs ved $u=x/x_0$, hvor $x_0$ er en referenceværdi med samme dimension. Hvis en model har formen $y=A\,f(x/x_0)$, beskriver $f$ en universel form, mens $A$ og $x_0$ sætter de fysiske skalaer.

For en positiv, differentiabel størrelse $y=f(x)$ med $x>0$ er den logaritmiske følsomhed
\[
\frac{\dd\ln y}{\dd\ln x}=\frac{x}{y}\frac{\dd y}{\dd x}.
\]
Den måler en relativ ændring i output pr. relativ ændring i input. For en potenslov $y=Cx^\alpha$ er følsomheden netop $\alpha$. Det hjælper med at afgøre, hvilke inputusikkerheder der er mest betydningsfulde, men usikkerhedsforplantning kræver yderligere antagelser om fejl og korrelationer.

\section{Eksponential- og logaritmefunktioner}
\subsection{Eksponentiel vækst og henfald}
For $a>0$ defineres $a^x=e^{x\ln a}$. Eksponentialfunktionen er positiv og opfylder
\[
e^{x+y}=e^x e^y,\qquad
\frac{\dd}{\dd x}e^{kx}=k e^{kx}.
\]
Differentialligningen $\dot y=ky$ har løsningen $y(t)=y(0)e^{kt}$ for konstant $k$. Her har $k$ dimension invers tid. Ved $k<0$ beskriver modellen eksponentielt henfald, og karakteristisk tid er $\tau=1/|k|$. Ved $k>0$ beskriver den eksponentiel vækst; fordoblingstiden er $\ln2/k$.

Logaritmen er invers til eksponentialfunktionen på de positive reelle tal. For $x,y>0$ og reel $\alpha$ gælder
\[
\ln(xy)=\ln x+\ln y,\qquad
\ln\frac{x}{y}=\ln x-\ln y,\qquad
\ln(x^\alpha)=\alpha\ln x.
\]
Man må ikke skrive $\ln(x+y)=\ln x+\ln y$. Logaritmen omdanner produkter til summer, ikke addition til en simpel operation.

% Original principskitse til algebra: lineær, eksponentiel og logaritmisk skalering.
\begin{figure}[htbp]
\centering
\begin{tikzpicture}
\begin{axis}[
 width=0.86\textwidth,height=6.3cm,
 axis lines=middle,xmin=-1.3,xmax=3.3,ymin=-1.4,ymax=4.6,
 samples=170,domain=-1.2:3.2,
 xlabel={$x$},ylabel={$y$},
 grid=major,grid style={gray!15},
 legend style={at={(0.98,0.03)},anchor=south east,fill=white,draw=none,font=\small}
]
\addplot[very thick,color=arkivblaa] {x};
\addlegendentry{$y=x$}
\addplot[very thick,color=arkivrod] {exp(x)};
\addlegendentry{$y=e^x$}
\addplot[very thick,color=black,domain=0.06:3.2] {ln(x)};
\addlegendentry{$y=\ln x$}
\addplot[only marks,mark=*,mark size=1.5pt,color=arkivrod] coordinates {(0,1)};
\addplot[only marks,mark=*,mark size=1.5pt,color=black] coordinates {(1,0)};
\end{axis}
\end{tikzpicture}
\caption{De tre forskellige former for afhængighed, vist med dimensionsløse variable: lineær vækst, eksponentialfunktion og logaritme. Eksponentialfunktionen og logaritmen er hinandens inverse efter spejling i linjen $y=x$. Grafen viser kun et endeligt udsnit af de respektive definitionsområder.}
\label{fig:algebra-eksplog}
\end{figure}


\subsection{Logaritmer, grene og fysiske enheder}
I en fysisk formel skal et logaritme- eller eksponentialargument være dimensionsløst. Udtrykket $\ln x$ er derfor ikke meningsfuldt, hvis $x$ f.eks. betegner en længde uden en angivet enhedsskala. Brug i stedet $\ln(x/x_0)$ for en positiv referencelængde $x_0$.

For $x>0$ gælder den nyttige relation
\[
\frac{\dd}{\dd x}\ln x=\frac{1}{x},
\]
og for $x\ne0$ er den reelle afledte af $\ln|x|$ også $1/x$. I kompleks analyse bliver logaritmen flerbladet, så en gren skal vælges særskilt; de almindelige reelle logaritmeregler må ikke overføres ukritisk til vilkårlige komplekse argumenter.

\section{Metodevalg, kontrol og anvendelser i fysik}
Når en fremtidig mekanik- eller elektromagnetismeopgave reduceres til algebra, kan følgende rækkefølge genbruges:
\begin{enumerate}
\item Bestem de ukendte størrelser, deres dimensioner og definitionsmængder.
\item Identificér mulige symmetrier, bevarede størrelser og dimensionsløse parametre, før du isolerer variable.
\item Vælg en algebraisk omskrivning, der bevarer ligningens løsninger; noter særskilt, når dette ikke er garanteret.
\item Brug en relevant repræsentation: faktorisering til rødder, fuldstændigt kvadrat til ekstrema, partialbrøker til integration eller logaritmer til produkter og skalering.
\item Indsæt resultatet i den oprindelige ligning, og kontrollér grænsetilfælde og enheder.
\end{enumerate}
For et central-kraftproblem bruges f.eks. løsning af en andengradsligning til radialvendepunkter, og fuldstændiggørelse af kvadratet kan hjælpe med at undersøge et potentiales form. I lineær mekanik kræver et stort koblet ligningssystem ofte matricemetoder snarere end manuel elimination; disse behandles i kapitlet om lineær algebra.

\paragraph{Opslag: hyppige fejl.}
Undgå division med mulige nulled, forkert fortegn ved uligheder, skjulte inversgrene, eksponentialfunktioner med dimensionsbærende argumenter og antagelser om, at en lokal linearisering gælder globalt. Hvis modellen introducerer en parameter, bør særlige værdier som nul og grænsetilfælde behandles eksplicit.

\chapter{Differentialregning}
\label{ch:diff}
Differentialregning beskriver lokal ændring. Den centrale idé er, at en tilstrækkeligt regulær funktion tæt på et punkt kan erstattes af en affin funktion med en kontrolleret fejl. Hældningen er samtidig en grænseværdi, en lineær afbildning og en lokal følsomhed. Disse tre synsvinkler forbinder de sædvanlige regneregler med fysikkens hastigheder, responskoefficienter og perturbative beskrivelser.

Kapitlet behandler funktioner af én reel variabel. Partielle afledte, Jacobimatricer og differentiabilitet i flere dimensioner hører til kapitel 8. De nævnes kun for at afgrænse, hvilke konklusioner der kan overføres.

\section{Grænseværdier, kontinuitet og lokal adfærd}
\label{sec:graense}
Lad $f$ være defineret i en punkteret omegn af $x_0$. Udsagnet $\lim_{x\to x_0}f(x)=L$ betyder, at der for hvert $\varepsilon>0$ findes et $\delta>0$, sådan at
\[
0<|x-x_0|<\delta\quad\Longrightarrow\quad |f(x)-L|<\varepsilon.
\]
Definitionen kræver ikke, at $f(x_0)$ findes. Kontinuitet i $x_0$ kræver derimod både en funktionsværdi og sammenfaldet $\lim_{x\to x_0}f(x)=f(x_0)$. En venstre- og en højregrænseværdi skal eksistere og være ens, før en endelig tosidet grænseværdi findes. Ved endepunkter arbejder man relativt til definitionsmængden.

Grænseværdier kan føres gennem summer og produkter, når de enkelte endelige grænseværdier findes. For kvotienter skal nævnerens grænseværdi desuden være forskellig fra nul. Ubestemte former som $0/0$ og $\infty-\infty$ er signaler om, at udtrykkets struktur skal undersøges; de er ikke regneregler.

For $h\to0$ betyder $r(h)=o(h)$, at $r(h)/h\to0$. Udsagnet $r(h)=O(h^2)$ betyder, at der findes $C,\delta>0$ med $|r(h)|\le C|h|^2$ for $0<|h|<\delta$. Lille-$o$ udtrykker en forsvindende relativ fejl, mens store-$O$ giver en størrelsesorden med en lokal konstant. Notationen siger ikke noget om en numerisk fejlgrænse, før konstanten er bestemt.

\section{Den afledte og differentialet}
\label{sec:def}
\begin{definition}
Lad $x_0$ være et indre punkt i intervallet $I$. Funktionen $f\colon I\to\R$ er differentiabel i $x_0$, hvis den endelige grænseværdi
\[
 f'(x_0)=\lim_{h\to0}\frac{f(x_0+h)-f(x_0)}{h}
\]
eksisterer. En uendelig hældning tæller ikke som differentiabilitet i denne definition.
\end{definition}
Sekanthældningen måler den gennemsnitlige ændring over et endeligt interval. Den afledte er grænsen af disse hældninger, når intervallet kollapser mod punktet. Den geometriske tangent har ligningen
\[
 y=f(x_0)+f'(x_0)(x-x_0).
\]
Den samme definition kan omformes til
\begin{equation}
 f(x_0+h)=f(x_0)+f'(x_0)h+r(h),\qquad r(h)=o(h).
 \label{eq:linear}
\end{equation}
Denne formulering er særlig vigtig: differentiabilitet er eksistensen af en lokal lineær respons med en fejl, der forsvinder hurtigere end tilvæksten. Differentialet er afbildningen $\dd f_{x_0}\colon h\mapsto f'(x_0)h$. Skrivemåden $\dd f=f'(x)\dd x$ udtrykker denne lineære sammenhæng; den faktiske ændring $\Delta f$ indeholder desuden restleddet.

Hvis $x$ har enhed $[x]$ og $f$ har enhed $[f]$, har $f'$ enhed $[f]/[x]$. Differentialet har samme enhed som $f$. En logaritme eller en eksponentialfunktion skal i en fysisk model have et dimensionsløst argument.

\begin{figure}[htbp]
\centering
\begin{tikzpicture}
\begin{axis}[width=.88\textwidth,height=65mm,axis lines=middle,xlabel={$x$},ylabel={$y$},xmin=-.2,xmax=2.3,ymin=-.8,ymax=4.8,legend pos=north west,legend style={font=\small},samples=100]
\addplot[arkivblaa,thick,domain=0:2.15]{x^2};\addlegendentry{$f(x)=x^2$}
\addplot[arkivrod,thick,domain=0:2.2]{2*x-1};\addlegendentry{Tangent ved $x_0=1$}
\addplot[black,dashed,domain=.4:2.15]{3*x-2};\addlegendentry{Sekant, $h=1$}
\fill[black] (axis cs:1,1) circle[radius=1.5pt];
\fill[black] (axis cs:2,4) circle[radius=1.5pt];
\end{axis}
\end{tikzpicture}
\caption{Sekanten gennem $(1,1)$ og $(2,4)$ har hældning 3; tangenten ved $(1,1)$ har hældning 2. For $h\to0$ går sekantens hældning $2+h$ mod 2. Akserne er dimensionsløse.}
\label{fig:sekant}
\end{figure}


\subsection{Kontinuitet er nødvendig, men ikke tilstrækkelig}
Af \eqref{eq:linear} følger $f(x_0+h)-f(x_0)\to0$, så differentiabilitet medfører kontinuitet. Den omvendte implikation gælder ikke. Funktionen $|x|$ er kontinuert i nul, men har ensidige afledte $-1$ og $1$ og dermed ingen tosidet afledt der.

En afledt behøver heller ikke være kontinuert. Funktionen $f(x)=x^2\sin(1/x)$ for $x\ne0$ og $f(0)=0$ har $f'(0)=0$, mens $f'(x)=2x\sin(1/x)-\cos(1/x)$ for $x\ne0$ ikke har en grænseværdi ved nul. Man må derfor ikke erstatte antagelsen $f\in C^1$ med blot ``$f$ differentiabel'', hvis et argument kræver kontinuitet af $f'$.

\section{Regneregler og sammensatte afhængigheder}
\label{sec:regler}
Nedenstående regler gælder i et punkt, hvor de nævnte funktioner er differentiable. Konstanterne $a,b$ er uafhængige af $x$. For kædereglen skal den indre funktion være differentiabel i $x$, og den ydre funktion i den tilsvarende værdi.

\begin{longtable}{@{}p{0.23\textwidth}p{0.40\textwidth}p{0.29\textwidth}@{}}
\caption{Regneregler. Forudsætninger er en del af formlerne.}\label{tab:regler}\\
\toprule Regel & Formel & Ekstra betingelse\\\midrule\endfirsthead
\toprule Regel & Formel & Ekstra betingelse\\\midrule\endhead
Linearitet & $(af+bg)'=af'+bg'$ & $a,b$ konstante\\[2mm]
Produkt & $(fg)'=f'g+fg'$ & Begge faktorer differentiable\\[2mm]
Kvotient & $\displaystyle (f/g)'=\frac{f'g-fg'}{g^2}$ & $g(x)\ne0$\\[3mm]
Kæderegel & $(f\circ g)'(x)=f'(g(x))g'(x)$ & Sammensætning defineret\\[2mm]
Reciprok & $(1/f)'=-f'/f^2$ & $f(x)\ne0$\\[2mm]
Logaritmisk & $(\ln|f|)'=f'/f$ & $f(x)\ne0$\\[2mm]
Potens af funktion & $(f^\alpha)'=\alpha f^{\alpha-1}f'$ & For reel $\alpha$ er $f>0$ en tilstrækkelig betingelse\\[2mm]
Variabel eksponent & $(f^g)'=f^g(g'\ln f+gf'/f)$ & $f>0$, begge differentiable\\
\bottomrule
\end{longtable}
Produktreglen kommer fra at tilføje og trække $f(x+h)g(x)$ i produktets differens. Efter division med $h$ giver kontinuiteten af $f$ den ene faktor og differentiabiliteten af $g$ den anden. Kædereglen udtrykker, at to lokale lineariseringer sammensættes: den ydre respons skal vurderes ved den indre funktions aktuelle værdi.

Ved tidsafhængighed $q=q(t)$ er $\dd f(q(t))/\dd t=f'(q(t))\dot q(t)$. Prikken betyder differentiation med hensyn til tid, ikke med hensyn til $q$. Denne skelnen forebygger, at en partiel afhængighed forveksles med en total tidsafledt i mekanik og termodynamik.

\subsection{Inverse og implicitte funktioner}
Hvis $f$ er $C^1$ nær $x_0$ og $f'(x_0)\ne0$, er $f$ lokalt invertibel, og
\[
 (f^{-1})'(y_0)=\frac{1}{f'(x_0)},\qquad y_0=f(x_0).
\]
Resultatet er lokalt. En funktion kan godt have flere adskilte monotone grene. Et punkt med $f'=0$ er ikke automatisk et punkt uden en invers, men den inverse kan miste differentiabilitet der.

For en relation $F(x,y)=0$ med $F\in C^1$ og $F_y(x_0,y_0)\ne0$ findes lokalt en funktion $y(x)$, og
\[
 \frac{\dd y}{\dd x}=-\frac{F_x}{F_y}.
\]
Her betegner $F_x,F_y$ partielle afledte, vurderet langs løsningen. Hvis $F_y=0$, er formlen ikke anvendelig; man må undersøge en anden lokal parametrisering eller en mulig singularitet.

\section{Standardafledte med definitionsområder}
\label{sec:standard}
Trigonometriske formler forudsætter radianmål. Ved gradmål kommer en faktor $\pi/180$ ind via kædereglen. De inverse trigonometriske funktioner forstås som deres sædvanlige reelle hovedgrene.

\begin{longtable}{@{}p{0.28\textwidth}p{0.32\textwidth}p{0.30\textwidth}@{}}
\caption{Elementære standardafledte på åbne reelle definitionsområder.}\\
\toprule $f(x)$ & $f'(x)$ & Område eller betingelse\\\midrule\endfirsthead
\toprule $f(x)$ & $f'(x)$ & Område eller betingelse\\\midrule\endhead
$c$ & $0$ & $c$ konstant\\[1mm]
$x^n$ & $nx^{n-1}$ & Heltal $n\ge1$; også nul, hvor udtrykket er defineret\\[1mm]
$x^\alpha$ & $\alpha x^{\alpha-1}$ & $x>0$, $\alpha\in\R$\\[1mm]
$\e^x$ & $\e^x$ & $x\in\R$\\[1mm]
$a^x$ & $a^x\ln a$ & $a>0$ konstant\\[1mm]
$\ln x$ & $1/x$ & $x>0$\\[1mm]
$\ln|x|$ & $1/x$ & $x\ne0$\\[1mm]
$\log_a x$ & $1/(x\ln a)$ & $x>0$, $a>0$, $a\ne1$\\[1mm]
$\sin x$ & $\cos x$ & $x\in\R$\\[1mm]
$\cos x$ & $-\sin x$ & $x\in\R$\\[1mm]
$\tan x$ & $1/\cos^2 x$ & $\cos x\ne0$\\[1mm]
$\cot x$ & $-1/\sin^2 x$ & $\sin x\ne0$\\[1mm]
$\arcsin x$ & $1/\sqrt{1-x^2}$ & $|x|<1$\\[1mm]
$\arccos x$ & $-1/\sqrt{1-x^2}$ & $|x|<1$\\[1mm]
$\arctan x$ & $1/(1+x^2)$ & $x\in\R$\\[1mm]
$\sinh x$ & $\cosh x$ & $x\in\R$\\[1mm]
$\cosh x$ & $\sinh x$ & $x\in\R$\\[1mm]
$\tanh x$ & $1/\cosh^2 x$ & $x\in\R$\\[1mm]
$\operatorname{arsinh}x$ & $1/\sqrt{1+x^2}$ & $x\in\R$\\[1mm]
$\operatorname{arcosh}x$ & $1/\sqrt{x^2-1}$ & $x>1$\\[1mm]
$\operatorname{artanh}x$ & $1/(1-x^2)$ & $|x|<1$\\
\bottomrule
\end{longtable}
En formel gælder kun, hvor den oprindelige funktion er defineret, og de anførte differentiabilitetskrav er opfyldt. Algebraisk forkortelse kan give et udtryk med et større definitionsområde end den oprindelige funktion; det udvider ikke automatisk denne funktion. For $|x|$ er den afledte $\operatorname{sgn}(x)$ for $x\ne0$, men ikke ved nul.

\section{Middelværdisætninger og globale konklusioner}
\label{sec:middel}
\begin{saetning}[Rolle]
Er $f$ kontinuert på $[a,b]$, differentiabel på $(a,b)$ og $f(a)=f(b)$, hvor $a<b$, findes et $c\in(a,b)$ med $f'(c)=0$.
\end{saetning}
Kontinuiteten på det kompakte interval sikrer, at maksimum og minimum antages. Hvis funktionen ikke er konstant, ligger mindst ét af disse ekstrema i det indre. Differentiabilitet tvinger da den afledte til at være nul der. Endepunktskontinuiteten og den indre differentiabilitet udfører altså forskellige opgaver i argumentet.

\begin{saetning}[Lagranges middelværdisætning]
Er $f$ kontinuert på $[a,b]$ og differentiabel på $(a,b)$, findes et $c\in(a,b)$, sådan at
\[
 f(b)-f(a)=f'(c)(b-a).
\]
\end{saetning}
Sætningen følger ved at trække sekantens affine funktion fra $f$ og anvende Rolle. Den siger, at den gennemsnitlige ændring svarer til en øjeblikkelig ændring et sted i intervallet. Den angiver generelt hverken et entydigt $c$ eller en konstruktiv metode til at finde det.

Hvis $|f'(x)|\le M$ på intervallet, følger stabilitetsestimatet
\[
 |f(b)-f(a)|\le M|b-a|.
\]
Det er en Lipschitzgrænse: en begrænset lokal følsomhed giver en kontrolleret endelig ændring. Hvis $f'\ge0$ på et interval, er $f$ ikke-aftagende. Hvis $f'>0$ overalt, er $f$ strengt voksende. Den omvendte slutning er stærkere end nødvendigt: en strengt voksende differentiabel funktion kan have $f'=0$ i enkelte punkter.

Cauchys middelværdisætning gælder for to funktioner med de samme kontinuitets- og differentiabilitetskrav:
\[
 [f(b)-f(a)]g'(c)=[g(b)-g(a)]f'(c)
\]
for et $c\in(a,b)$. Man må først dividere med $g'(c)$ og $g(b)-g(a)$, når de er forskellige fra nul. Denne sætning ligger bag l'Hôpitals regel.

\section{Højere afledte, monotoni og krumning}
\label{sec:hoejere}
Den anden afledte $f''$ beskriver ændringen i hældningen. Klassen $C^n(I)$ består af funktioner med kontinuerte afledte til og med orden $n$ på $I$. $C^\infty$ betyder glat, men betyder ikke nødvendigvis analytisk: en funktion kan være uendeligt differentiabel uden at være lig sin Taylorrække.

\subsection{Ekstrema og stationære punkter}
For en differentiabel funktion er $f'(x_0)=0$ en nødvendig betingelse for et indre lokalt ekstremum. Den er ikke tilstrækkelig. Fortegnsskiftet af $f'$ afgør typen: fra negativ til positiv giver et minimum; fra positiv til negativ giver et maksimum. Intet fortegnsskift giver normalt intet ekstremum, men degenererede tilfælde kræver en direkte undersøgelse.

Hvis $f\in C^2$ nær $x_0$, $f'(x_0)=0$ og $f''(x_0)>0$, er punktet et strengt lokalt minimum; ved $f''(x_0)<0$ et strengt lokalt maksimum. Hvis $f''(x_0)=0$, er testen uafgjort. Mere generelt: er den første ikke-nul afledte efter konstantleddet af orden $m$, og gælder en Taylorudvikling med restled $o((x-x_0)^m)$, giver lige $m$ et ekstremum bestemt af fortegnet; ulige $m$ giver intet ekstremum.

En fuldstændig undersøgelse på et lukket interval skal også omfatte endepunkter og indre punkter, hvor funktionen ikke er differentiabel. En liste over løsninger til $f'=0$ alene er ikke en optimeringsmetode på hele definitionsmængden.

\subsection{Konveksitet og geometrisk krumning}
En differentiabel funktion er konveks på et interval, når dens graf ligger over alle sine tangenter. For $C^2$-funktioner er dette ækvivalent med $f''\ge0$. Konkavitet svarer til $f''\le0$. Et vendepunkt kræver skift mellem konveksitet og konkavitet; ligningen $f''=0$ alene er utilstrækkelig.

For en plan graf $y=f(x)$ med dimensionsløse, lige skalerede koordinater er den geometriske krumnings størrelse
\[
 \kappa(x)=\frac{|f''(x)|}{(1+f'(x)^2)^{3/2}}.
\]
Andenafledte og geometrisk krumning er altså forskellige størrelser. Hvis akserne har forskellige fysiske enheder, skal en fysisk indlejring og skalering fastlægges, før denne geometriske formel kan fortolkes.

\begin{figure}[htbp]
\centering
\begin{tikzpicture}
\begin{axis}[width=.88\textwidth,height=65mm,axis lines=middle,xlabel={$x$},ylabel={$y$},xmin=-2,xmax=2,ymin=-3.5,ymax=3.5,legend pos=north west,legend style={font=\small},samples=150]
\addplot[arkivblaa,thick,domain=-1.9:1.9]{x^3-3*x};\addlegendentry{$f=x^3-3x$}
\addplot[arkivrod,dashed,domain=-1.45:1.45]{3*x^2-3};\addlegendentry{$f'=3x^2-3$}
\fill[black] (axis cs:-1,2) circle[radius=1.5pt];
\fill[black] (axis cs:1,-2) circle[radius=1.5pt];
\fill[black] (axis cs:0,0) circle[radius=1.5pt];
\end{axis}
\end{tikzpicture}
\caption{Afledtens fortegn styrer monotonien. Ved $x=-1$ ligger et maksimum, ved $x=1$ et minimum. Vendepunktet $x=0$ har $f''=0$ med fortegnsskift, men er ikke stationært. Akserne er dimensionsløse.}
\end{figure}


\section{Taylorudvikling og kontrollerede restled}
\label{sec:taylor}
Taylorpolynomiet af grad $n$ omkring $a$ er
\[
 T_n(x)=\sum_{k=0}^n\frac{f^{(k)}(a)}{k!}(x-a)^k.
\]
Det matcher funktionsværdien og de første $n$ afledte i udviklingspunktet. Polynomiet beskriver lokal struktur; det er ikke en global erstatning uden en fejlundersøgelse.

\begin{saetning}[Taylor med Lagrangerest]
Lad $f\in C^{n+1}$ på et åbent interval, der indeholder segmentet mellem $a$ og $x$. Da findes et $\xi$ mellem $a$ og $x$, sådan at
\begin{equation}
 f(x)=T_n(x)+\frac{f^{(n+1)}(\xi)}{(n+1)!}(x-a)^{n+1}.
 \label{eq:rest}
\end{equation}
\end{saetning}
Den valgte $C^{n+1}$-antagelse er en enkel tilstrækkelig betingelse. Hvis $|f^{(n+1)}(t)|\le M$ på segmentet, fås den praktiske fejlgrænse
\[
 |f(x)-T_n(x)|\le\frac{M|x-a|^{n+1}}{(n+1)!}.
\]
Ved blot $f\in C^n$ nær $a$ findes Peanoformen $f(x)=T_n(x)+o((x-a)^n)$ for $x\to a$. Den form giver ikke i sig selv en kendt konstant for fejlens størrelse. Når $f^{(n+1)}$ er kontinuert, er også integralrestleddet gyldigt:
\[
 R_n(x)=\frac{1}{n!}\int_a^x (x-t)^n f^{(n+1)}(t)\dd t.
\]
Integralformen er velegnet til fortegns- og normestimater og gælder også for $x<a$ med integralets sædvanlige orientering.

\subsection{Standardudviklinger omkring nul}
Her er $z$ dimensionsløs, og restleddene gælder for $z\to0$:
\begin{align*}
 \e^z&=1+z+\tfrac12z^2+\tfrac16z^3+O(z^4),\\
 \sin z&=z-\tfrac16z^3+O(z^5),\\
 \cos z&=1-\tfrac12z^2+\tfrac1{24}z^4+O(z^6),\\
 \ln(1+z)&=z-\tfrac12z^2+\tfrac13z^3+O(z^4),\\
 (1+z)^\alpha&=1+\alpha z+\tfrac12\alpha(\alpha-1)z^2+O(z^3).
\end{align*}
Logaritmen og den reelle potens kræver $1+z>0$. De tilhørende uendelige potensrækker har yderligere konvergensbetingelser: logaritme- og den generelle binomialrække har konvergensradius 1, med særskilt analyse på randen; for ikke-negative heltalseksponenter stopper binomialrækken og er et polynomium. Eksponential-, sinus- og cosinusrækkerne konvergerer for alle reelle $z$.

At alle afledte findes, sikrer ikke, at en uendelig Taylorrække repræsenterer funktionen. Funktionen $\e^{-1/x^2}$ for $x\ne0$, udvidet med værdien nul i nul, er glat og har alle afledte lig nul i nul, men er positiv uden for nul. Analyticitet er en ekstra egenskab.

\begin{figure}[htbp]
\centering
\begin{tikzpicture}
\begin{axis}[width=.88\textwidth,height=65mm,axis lines=middle,xlabel={$x$},ylabel={$y$},xmin=-2,xmax=2,ymin=-1.2,ymax=1.2,legend pos=north west,legend style={font=\small},samples=150]
\addplot[arkivblaa,thick,domain=-2:2]{sin(deg(x))};\addlegendentry{$\sin x$}
\addplot[black,dotted,domain=-1.2:1.2]{x};\addlegendentry{$T_1=x$}
\addplot[arkivrod,dashed,thick,domain=-2:2]{x-x^3/6};\addlegendentry{$T_3=x-x^3/6$}
\end{axis}
\end{tikzpicture}
\caption{Taylorpolynomier omkring nul. Den kubiske tilnærmelse følger sinus tættere nær nul end tangenten. For alle reelle $x$ giver Taylor til orden 4, hvor koefficienten til $x^4$ er nul, grænsen $|\sin x-T_3(x)|\le |x|^5/120$. En lille relativ fejl følger ikke automatisk nær sinusfunktionens nulpunkter.}
\end{figure}


\section{Grænseværdier og l'Hôpitals regel}
\label{sec:hopital}
En anvendelig version er følgende: Lad $f,g$ være differentiable på et punkteret ensidigt interval mod et endepunkt $a$, med $g'\ne0$. Antag enten $f,g\to0$ eller $|f|,|g|\to\infty$. Hvis $f'/g'$ har en grænseværdi $L$, endelig eller uendelig, gælder under disse betingelser
\[
 \lim_{x\to a}\frac{f(x)}{g(x)}=L.
\]
Samme princip kan anvendes ved uendelige endepunkter. Tosidige grænseværdier kræver kontrol fra begge sider. Ved gentagen brug skal betingelserne verificeres igen på hvert trin. Hvis kvotienten af de afledte ikke har en grænseværdi, er testen uafgjort; det beviser ikke, at den oprindelige kvotient mangler en grænseværdi.

Reglen differentierer tæller og nævner hver for sig til en ny kvotient; den beregner ikke den afledte af den oprindelige kvotient. Former som $0\cdot\infty$ skal først omskrives til en relevant kvotient. Taylorudvikling er ofte mere informativ, fordi den samtidig viser den første ikke-nul orden og tilnærmelsens fejl.

\section{Fysisk fortolkning og følsomhed}
\label{sec:fysik}
For en tilstrækkeligt glat position $q(t)$ er hastigheden $v=\dot q$ og accelerationen $a=\ddot q$. De har enheder henholdsvis længde pr. tid og længde pr. tid i anden. En grafisk hældning afhænger af aksernes skala, mens den fysiske afledte bestemmes af variablerne og deres enheder.

For en éndimensional konservativ kraft er $F(q)=-U'(q)$. Et ligevægtspunkt opfylder $U'(q_0)=0$. Hvis $U''(q_0)>0$ og massen $m>0$, giver lineariseringen
\[
 m\ddot\eta=-U''(q_0)\eta+O(\eta^2),\qquad \eta=q-q_0,
\]
forudsat tilstrækkelig regularitet, eksempelvis $U\in C^3$ nær $q_0$. Den lineære svingningsfrekvens er $\omega=\sqrt{U''(q_0)/m}$. For $U''(q_0)<0$ er den lineariserede ligevægt ustabil. Ved $U''(q_0)=0$ må højere ordener undersøges. Dette er en generel kobling mellem afledte og stabilitet, ikke en konkret opgaveløsning.

For en målt variabel $x$ giver en lille perturbation $\delta x$ den førsteordensændring $\delta f\approx f'(x)\delta x$. Hvis $x\ne0$ og $f(x)\ne0$, er den relative følsomhed
\[
 \kappa_{\mathrm{rel}}(x)=\left|\frac{x f'(x)}{f(x)}\right|,\qquad
 \left|\frac{\delta f}{f}\right|\approx\kappa_{\mathrm{rel}}\left|\frac{\delta x}{x}\right|.
\]
Nær nulpunkter kan relativ fejl være uegnet, selv når den absolutte fejl er lille. For en stokastisk usikkerhed med lille standardafvigelse $\sigma_x$ giver lineariseringen $\sigma_f\approx|f'(x)|\sigma_x$. Ikke-linearitet, skæve fordelinger og korrelationer kræver en videre analyse i statistikkapitlet; formlen er ikke en generel eksakt usikkerhedslov.

\section{Numerisk differentiation og fejlbalance}
\label{sec:numerisk}
Når $f$ kun er tilgængelig gennem evalueringer, kan differenskvotienter approksimere den afledte. For $h\ne0$ er de mest brugte:
\begin{align*}
 D_h^+f(x)&=\frac{f(x+h)-f(x)}{h}
           =f'(x)+\tfrac h2f''(x)+O(h^2),\\
 D_h^0f(x)&=\frac{f(x+h)-f(x-h)}{2h}
           =f'(x)+\tfrac{h^2}{6}f^{(3)}(x)+O(h^4),\\
 D_h^2f(x)&=\frac{f(x+h)-2f(x)+f(x-h)}{h^2}
           =f''(x)+\tfrac{h^2}{12}f^{(4)}(x)+O(h^4).
\end{align*}
Tilstrækkelige lokale regularitetskrav for de viste restled er henholdsvis $C^3$, $C^5$ og $C^6$. De centrerede formler udnytter symmetri til at eliminere ulige eller lige bidrag. Ved rande kræver de funktionsværdier uden for intervallet; brug i stedet en passende ensidig formel eller en fysisk begrundet randudvidelse.

Mindre $h$ er ikke altid bedre. Hvis hver funktionsværdi har absolut fejl højst $\varepsilon_f$, kan fejlbidraget til den centrerede førsteafledte være op til $\varepsilon_f/|h|$. En simpel model for totalfejlen er derfor
\[
 E(h)\lesssim C h^2+\frac{\varepsilon_f}{|h|}.
\]
Her kommer første led fra trunkering og andet led fra afrunding eller datastøj. Balancen giver skalaen $|h|\sim(\varepsilon_f/C)^{1/3}$ op til konstante faktorer, når $C>0$. Dette er et lokalt skalaestimat, ikke et universelt optimalt skridt. Variabelskalering, støjstruktur og ukendte højere afledte påvirker valget.

\begin{figure}[htbp]
\centering
\begin{tikzpicture}
\begin{loglogaxis}[width=.88\textwidth,height=65mm,xlabel={Skridt $h$},ylabel={Fejlmodel},xmin=1e-7,xmax=1e-1,ymin=1e-11,ymax=1e-1,legend pos=north east,legend style={font=\small},samples=100]
\addplot[arkivblaa,dashed,domain=1e-7:1e-1]{x^2};\addlegendentry{$h^2$}
\addplot[arkivrod,dotted,thick,domain=1e-7:1e-1]{1e-12/x};\addlegendentry{$10^{-12}/h$}
\addplot[black,thick,domain=1e-7:1e-1]{x^2+1e-12/x};\addlegendentry{Sum}
\end{loglogaxis}
\end{tikzpicture}
\caption{Illustrativ fejlmodel for centreret differentiation med $C=1$ og $\varepsilon_f=10^{-12}$ i dimensionsløse variable. Minimum ligger ved $h=(\varepsilon_f/2C)^{1/3}\approx7{,}94\cdot10^{-5}$. Kurverne er en analytisk model, ikke målte flydende-kommafejl.}
\end{figure}


Automatisk differentiation beregner afledte ved at føre kædereglen gennem programmets operationer. Den undgår differenskvotientens trunkeringsfejl, men hverken diskontinuerte programgrene, afrunding eller en forkert fysisk model forsvinder af den grund. Symbolsk differentiation, automatisk differentiation og numeriske differenser er forskellige redskaber.

\clearpage
\section{Metodeoverblik og kontrolpunkter}
\label{sec:kontrol}
Begynd med definitionsmængde og ønsket konklusion. Til lokal linearisering bruges differentiabilitet; til en kvantitativ fejlgrænse kræves kontrol over relevante højere afledte. Til global monotoni anvendes fortegnet på den afledte på hele det relevante interval. Til ekstremumskontrol undersøges både stationære punkter, singulariteter og rande.

Ved sammensætninger skal den indre funktions afledte og den ydre funktions evalueringspunkt fremgå. Ved inverse funktioner angives gren og ikke-nul hældning. Ved Taylorapproksimation angives udviklingspunkt, orden, gyldighedsområde og restled. Ved numerisk differentiation afprøves flere skridtstørrelser, så både konvergensområdet og støjgulvet kan identificeres.

En faglig kontrol omfatter desuden:
\begin{itemize}
\item Har alle symboler en definition, og er variable og konstanter skilt ad?
\item Stemmer enhederne, også efter tidsdifferentiation og kædereglen?
\item Er udsagnet lokalt eller globalt, og ligger punktet i det indre eller på randen?
\item Er der divideret med en størrelse, som kan være nul?
\item Er en nødvendig betingelse fejlagtigt brugt som tilstrækkelig?
\item Matcher figurernes akser, kurver og markeringer de beskrevne formler?
\end{itemize}
Disse kontroller er en del af den matematiske argumentation. De erstatter ikke et bevis, men gør skjulte antagelser synlige og forbinder håndbogens formler med forsvarlig anvendelse i fysik.

\chapter{Integralregning og standardintegraler}
\planlagt{Definitioner, notation, forudsætninger, centrale resultater og fysikforbindelser udvikles samlet.}
\section{Stamfunktioner og analysens hovedsætning}
\leavevmode\par
% Redaktionelt: tilføj definitioner, gyldighed, relationer og relevante figurer.

\section{Substitution, partiel integration og partialbrøker}
\leavevmode\par
% Redaktionelt: tilføj definitioner, gyldighed, relationer og relevante figurer.

\section{Standardintegraler med gyldighedsbetingelser}
\leavevmode\par
% Redaktionelt: tilføj definitioner, gyldighed, relationer og relevante figurer.

\section{Uegentlige integraler og parameterintegration}
\leavevmode\par
% Redaktionelt: tilføj definitioner, gyldighed, relationer og relevante figurer.

\chapter{Rækker, konvergens og approksimationer}
\planlagt{Definitioner, notation, forudsætninger, centrale resultater og fysikforbindelser udvikles samlet.}
\section{Følger, punktvis og uniform konvergens}
\leavevmode\par
% Redaktionelt: tilføj definitioner, gyldighed, relationer og relevante figurer.

\section{Konvergenskriterier og potensrækker}
\leavevmode\par
% Redaktionelt: tilføj definitioner, gyldighed, relationer og relevante figurer.

\section{Taylor- og Fourierrækker}
\leavevmode\par
% Redaktionelt: tilføj definitioner, gyldighed, relationer og relevante figurer.

\section{Asymptotiske rækker og fejlkontrol}
\leavevmode\par
% Redaktionelt: tilføj definitioner, gyldighed, relationer og relevante figurer.

\chapter{Differentialligninger}
\planlagt{Definitioner, notation, forudsætninger, centrale resultater og fysikforbindelser udvikles samlet.}
\section{Eksistens, entydighed og begyndelsesbetingelser}
\leavevmode\par
% Redaktionelt: tilføj definitioner, gyldighed, relationer og relevante figurer.

\section{Førsteordensligninger og integrerende faktor}
\leavevmode\par
% Redaktionelt: tilføj definitioner, gyldighed, relationer og relevante figurer.

\section{Lineære systemer og fundamentalløsninger}
\leavevmode\par
% Redaktionelt: tilføj definitioner, gyldighed, relationer og relevante figurer.

\section{Stabilitet, faseportrætter og perturbation}
\leavevmode\par
% Redaktionelt: tilføj definitioner, gyldighed, relationer og relevante figurer.

\chapter{Partielle differentialligninger}
\planlagt{Definitioner, notation, forudsætninger, centrale resultater og fysikforbindelser udvikles samlet.}
\section{Klassifikation og randbetingelser}
\leavevmode\par
% Redaktionelt: tilføj definitioner, gyldighed, relationer og relevante figurer.

\section{Separation af variable}
\leavevmode\par
% Redaktionelt: tilføj definitioner, gyldighed, relationer og relevante figurer.

\section{Varme-, bølge- og Laplaceligningen}
\leavevmode\par
% Redaktionelt: tilføj definitioner, gyldighed, relationer og relevante figurer.

\section{Karakteristikker og svage løsninger}
\leavevmode\par
% Redaktionelt: tilføj definitioner, gyldighed, relationer og relevante figurer.

\chapter{Lineær algebra}
\planlagt{Definitioner, notation, forudsætninger, centrale resultater og fysikforbindelser udvikles samlet.}
\section{Vektorrum, basis og lineære afbildninger}
\leavevmode\par
% Redaktionelt: tilføj definitioner, gyldighed, relationer og relevante figurer.

\section{Determinanter, rang og lineære ligninger}
\leavevmode\par
% Redaktionelt: tilføj definitioner, gyldighed, relationer og relevante figurer.

\section{Egenværdier og diagonalisering}
\leavevmode\par
% Redaktionelt: tilføj definitioner, gyldighed, relationer og relevante figurer.

\section{Indre produkter, spektralsætning og SVD}
\leavevmode\par
% Redaktionelt: tilføj definitioner, gyldighed, relationer og relevante figurer.

\chapter{Flervariabel calculus}
\planlagt{Definitioner, notation, forudsætninger, centrale resultater og fysikforbindelser udvikles samlet.}
\section{Differentiabilitet, Jacobimatrix og kæderegel}
\leavevmode\par
% Redaktionelt: tilføj definitioner, gyldighed, relationer og relevante figurer.

\section{Hessian og kritiske punkter}
\leavevmode\par
% Redaktionelt: tilføj definitioner, gyldighed, relationer og relevante figurer.

\section{Implicitte funktioner og Lagrangemultiplikatorer}
\leavevmode\par
% Redaktionelt: tilføj definitioner, gyldighed, relationer og relevante figurer.

\section{Integration i krumme koordinater og Jacobideterminant}
\leavevmode\par
% Redaktionelt: tilføj definitioner, gyldighed, relationer og relevante figurer.

\chapter{Vektorregning}
\planlagt{Definitioner, notation, forudsætninger, centrale resultater og fysikforbindelser udvikles samlet.}
\section{Gradient, divergens og rotation}
\leavevmode\par
% Redaktionelt: tilføj definitioner, gyldighed, relationer og relevante figurer.

\section{Kurve- og fladeintegraler med orientering}
\leavevmode\par
% Redaktionelt: tilføj definitioner, gyldighed, relationer og relevante figurer.

\section{Gauss og Stokes med forudsætninger}
\leavevmode\par
% Redaktionelt: tilføj definitioner, gyldighed, relationer og relevante figurer.

\section{Ortogonal krum koordinatnotation}
\leavevmode\par
% Redaktionelt: tilføj definitioner, gyldighed, relationer og relevante figurer.

\chapter{Analytisk geometri i 2D og 3D}
\planlagt{Definitioner, notation, forudsætninger, centrale resultater og fysikforbindelser udvikles samlet.}
\section{Linjer, planer og afstande}
\leavevmode\par
% Redaktionelt: tilføj definitioner, gyldighed, relationer og relevante figurer.

\section{Keglesnit og kvadrikker}
\leavevmode\par
% Redaktionelt: tilføj definitioner, gyldighed, relationer og relevante figurer.

\section{Parametriserede kurver, krumning og torsion}
\leavevmode\par
% Redaktionelt: tilføj definitioner, gyldighed, relationer og relevante figurer.

\section{Koordinatsystemer og orientering}
\leavevmode\par
% Redaktionelt: tilføj definitioner, gyldighed, relationer og relevante figurer.

\chapter{Komplekse tal og kompleks analyse}
\planlagt{Definitioner, notation, forudsætninger, centrale resultater og fysikforbindelser udvikles samlet.}
\section{Komplekse tal, argument og grene}
\leavevmode\par
% Redaktionelt: tilføj definitioner, gyldighed, relationer og relevante figurer.

\section{Holomorfi og Cauchy–Riemann}
\leavevmode\par
% Redaktionelt: tilføj definitioner, gyldighed, relationer og relevante figurer.

\section{Konturintegration og Cauchys sætning}
\leavevmode\par
% Redaktionelt: tilføj definitioner, gyldighed, relationer og relevante figurer.

\section{Laurentrækker, singulariteter og residuer}
\leavevmode\par
% Redaktionelt: tilføj definitioner, gyldighed, relationer og relevante figurer.

\chapter{Fourier- og Laplacetransformationer}
\planlagt{Definitioner, notation, forudsætninger, centrale resultater og fysikforbindelser udvikles samlet.}
\section{Konventioner, normalisering og enheder}
\leavevmode\par
% Redaktionelt: tilføj definitioner, gyldighed, relationer og relevante figurer.

\section{Fouriertransformation og inversion}
\leavevmode\par
% Redaktionelt: tilføj definitioner, gyldighed, relationer og relevante figurer.

\section{Laplace og konvergensområder}
\leavevmode\par
% Redaktionelt: tilføj definitioner, gyldighed, relationer og relevante figurer.

\section{Parseval, spektra og differentialoperatorer}
\leavevmode\par
% Redaktionelt: tilføj definitioner, gyldighed, relationer og relevante figurer.

\chapter{Konvolution, deltafunktion og Greenfunktioner}
\planlagt{Definitioner, notation, forudsætninger, centrale resultater og fysikforbindelser udvikles samlet.}
\section{Testfunktioner og distributioner}
\leavevmode\par
% Redaktionelt: tilføj definitioner, gyldighed, relationer og relevante figurer.

\section{Deltafunktionen og koordinatskift}
\leavevmode\par
% Redaktionelt: tilføj definitioner, gyldighed, relationer og relevante figurer.

\section{Konvolution og fundamentalløsninger}
\leavevmode\par
% Redaktionelt: tilføj definitioner, gyldighed, relationer og relevante figurer.

\section{Greenfunktioner, randbetingelser og kausalitet}
\leavevmode\par
% Redaktionelt: tilføj definitioner, gyldighed, relationer og relevante figurer.

\chapter{Sturm–Liouville-teori}
\planlagt{Definitioner, notation, forudsætninger, centrale resultater og fysikforbindelser udvikles samlet.}
\section{Selvadjungerede operatorer og vægtfunktion}
\leavevmode\par
% Redaktionelt: tilføj definitioner, gyldighed, relationer og relevante figurer.

\section{Randbetingelser og randled}
\leavevmode\par
% Redaktionelt: tilføj definitioner, gyldighed, relationer og relevante figurer.

\section{Ortogonalitet, spektrum og fuldstændighed}
\leavevmode\par
% Redaktionelt: tilføj definitioner, gyldighed, relationer og relevante figurer.

\section{Egenfunktionsudvikling og Greenkerner}
\leavevmode\par
% Redaktionelt: tilføj definitioner, gyldighed, relationer og relevante figurer.

\chapter{Specialfunktioner}
\planlagt{Definitioner, notation, forudsætninger, centrale resultater og fysikforbindelser udvikles samlet.}
\section{Gamma- og betafunktionen}
\leavevmode\par
% Redaktionelt: tilføj definitioner, gyldighed, relationer og relevante figurer.

\section{Besselfunktioner og cylindrisk symmetri}
\leavevmode\par
% Redaktionelt: tilføj definitioner, gyldighed, relationer og relevante figurer.

\section{Legendre og sfæriske harmoniske}
\leavevmode\par
% Redaktionelt: tilføj definitioner, gyldighed, relationer og relevante figurer.

\section{Hermite, Laguerre og ortogonalitet}
\leavevmode\par
% Redaktionelt: tilføj definitioner, gyldighed, relationer og relevante figurer.

\chapter{Sandsynlighed, statistik og usikkerheder}
\planlagt{Definitioner, notation, forudsætninger, centrale resultater og fysikforbindelser udvikles samlet.}
\section{Fordelinger, momenter og betingning}
\leavevmode\par
% Redaktionelt: tilføj definitioner, gyldighed, relationer og relevante figurer.

\section{Estimation, likelihood og konfidensintervaller}
\leavevmode\par
% Redaktionelt: tilføj definitioner, gyldighed, relationer og relevante figurer.

\section{Usikkerhedspropagation og korrelation}
\leavevmode\par
% Redaktionelt: tilføj definitioner, gyldighed, relationer og relevante figurer.

\section{Regression og modelkontrol}
\leavevmode\par
% Redaktionelt: tilføj definitioner, gyldighed, relationer og relevante figurer.

\chapter{Numeriske metoder}
\planlagt{Definitioner, notation, forudsætninger, centrale resultater og fysikforbindelser udvikles samlet.}
\section{Afrunding, kondition og stabilitet}
\leavevmode\par
% Redaktionelt: tilføj definitioner, gyldighed, relationer og relevante figurer.

\section{Rodfinding og lineære systemer}
\leavevmode\par
% Redaktionelt: tilføj definitioner, gyldighed, relationer og relevante figurer.

\section{Interpolation, differentiation og kvadratur}
\leavevmode\par
% Redaktionelt: tilføj definitioner, gyldighed, relationer og relevante figurer.

\section{ODE-integration, konvergens og fejlestimation}
\leavevmode\par
% Redaktionelt: tilføj definitioner, gyldighed, relationer og relevante figurer.

\chapter{Variationsregning}
\planlagt{Definitioner, notation, forudsætninger, centrale resultater og fysikforbindelser udvikles samlet.}
\section{Funktionaler og første variation}
\leavevmode\par
% Redaktionelt: tilføj definitioner, gyldighed, relationer og relevante figurer.

\section{Euler–Lagrange og randbetingelser}
\leavevmode\par
% Redaktionelt: tilføj definitioner, gyldighed, relationer og relevante figurer.

\section{Bivariable, bindinger og multiplikatorer}
\leavevmode\par
% Redaktionelt: tilføj definitioner, gyldighed, relationer og relevante figurer.

\section{Symmetrier og Noethers sætning}
\leavevmode\par
% Redaktionelt: tilføj definitioner, gyldighed, relationer og relevante figurer.

\chapter{Gruppeteori og symmetrier}
\planlagt{Definitioner, notation, forudsætninger, centrale resultater og fysikforbindelser udvikles samlet.}
\section{Grupper, homomorfier og repræsentationer}
\leavevmode\par
% Redaktionelt: tilføj definitioner, gyldighed, relationer og relevante figurer.

\section{Liegrupper, Liealgebraer og generatorer}
\leavevmode\par
% Redaktionelt: tilføj definitioner, gyldighed, relationer og relevante figurer.

\section{Rotationer, SO(3) og SU(2)}
\leavevmode\par
% Redaktionelt: tilføj definitioner, gyldighed, relationer og relevante figurer.

\section{Symmetrier og bevarelseslove}
\leavevmode\par
% Redaktionelt: tilføj definitioner, gyldighed, relationer og relevante figurer.

\chapter{Tensorregning}
\planlagt{Definitioner, notation, forudsætninger, centrale resultater og fysikforbindelser udvikles samlet.}
\section{Multilineære afbildninger og dualrum}
\leavevmode\par
% Redaktionelt: tilføj definitioner, gyldighed, relationer og relevante figurer.

\section{Indekser og transformationslove}
\leavevmode\par
% Redaktionelt: tilføj definitioner, gyldighed, relationer og relevante figurer.

\section{Kontraktion, symmetri og tensorprodukter}
\leavevmode\par
% Redaktionelt: tilføj definitioner, gyldighed, relationer og relevante figurer.

\section{Metrik, indeksflytning og orientering}
\leavevmode\par
% Redaktionelt: tilføj definitioner, gyldighed, relationer og relevante figurer.

\chapter{Differentialgeometri og matematik til GR}
\planlagt{Definitioner, notation, forudsætninger, centrale resultater og fysikforbindelser udvikles samlet.}
\section{Mangfoldigheder og koordinatkort}
\leavevmode\par
% Redaktionelt: tilføj definitioner, gyldighed, relationer og relevante figurer.

\section{Tangentrum og dualrum}
\leavevmode\par
% Redaktionelt: tilføj definitioner, gyldighed, relationer og relevante figurer.

\section{Tensorer og koordinattransformationer}
\leavevmode\par
% Redaktionelt: tilføj definitioner, gyldighed, relationer og relevante figurer.

\section{Metrikker og linjeelementer}
\leavevmode\par
% Redaktionelt: tilføj definitioner, gyldighed, relationer og relevante figurer.

\section{Kovariant differentiation og Christoffelsymboler}
\leavevmode\par
% Redaktionelt: tilføj definitioner, gyldighed, relationer og relevante figurer.

\section{Paralleltransport og holonomi}
\leavevmode\par
% Redaktionelt: tilføj definitioner, gyldighed, relationer og relevante figurer.

\section{Geodæter og affin parameter}
\leavevmode\par
% Redaktionelt: tilføj definitioner, gyldighed, relationer og relevante figurer.

\section{Riemanntensor og fortegnskonvention}
\leavevmode\par
% Redaktionelt: tilføj definitioner, gyldighed, relationer og relevante figurer.

\section{Riccitensor og skalarkrumning}
\leavevmode\par
% Redaktionelt: tilføj definitioner, gyldighed, relationer og relevante figurer.

\section{Geodætisk deviation}
\leavevmode\par
% Redaktionelt: tilføj definitioner, gyldighed, relationer og relevante figurer.

\section{Killingvektorer og Lie-afledte}
\leavevmode\par
% Redaktionelt: tilføj definitioner, gyldighed, relationer og relevante figurer.

\section{Differentialformer og Hodge-dualitet}
\leavevmode\par
% Redaktionelt: tilføj definitioner, gyldighed, relationer og relevante figurer.

\section{Bianchi-identiteterne}
\leavevmode\par
% Redaktionelt: tilføj definitioner, gyldighed, relationer og relevante figurer.

\section{Lorentzgeometri og introduktion til GR}
\leavevmode\par
% Redaktionelt: tilføj definitioner, gyldighed, relationer og relevante figurer.

\backmatter
\chapter{Kilder og faglig kontrol}
Pilotkapitlet er selvstændigt skrevet. Følgende standardværker er relevante til videre faglig kontrol, men der angives ingen verificerede sidetal eller kursusbindinger:
\begin{itemize}
\item Tom M. Apostol: \emph{Calculus}, bind 1. Differentialregning og Taylors sætning.
\item Michael Spivak: \emph{Calculus}. Grænseværdier, differentiabilitet og middelværdisætninger.
\item Bengt Fornberg: \emph{A Practical Guide to Pseudospectral Methods}. Numerisk differentiation og approksimation.
\end{itemize}
Konkret pensum og brugerens originale noter registreres først, når de er tilgængelige. Kontrolstatus og udviklingsprocedure fremgår af repositoryets dokumentation.
\end{document}
